\chapter{Command Reference}
\label{ch:cmdref}

Every command on every menu, in alphabetical order. Each entry says what the
command does, how to run it, and what it will not do.

Commands are named here as they appear on the menu. \menupath{Data > Sort up}
is listed under \menupath{Sort up}, and the chart commands under their own
names --- \menupath{Bars}, \menupath{Line} and \menupath{Pie}.

\section{The commands at a glance}

\begin{cctablethree}{Command}{Key}{Menu}{1.35in}{0.85in}{2.6in}
New & \key{F1} & File \\
Open & \key{F3} & File \\
Save & \key{F4} & File \\
Save as & \key{F6} & File \\
Import & \key{F7} & File \\
Export & \key{F8} & File \\
Quit & \key{F9} & File \\
Edit cell & \key{F2} & Edit \\
Select & \keyc{Ctrl}{A} & Edit \\
Copy & \keyc{Ctrl}{C} & Edit \\
Paste & \keyc{Ctrl}{V} & Edit \\
Undo & \key{F5} & Edit \\
Clear & \key{Delete} & Edit \\
Find & \keyc{Ctrl}{F} & Edit \\
Replace\ccdots & --- & Edit \\
Insert row & --- & Layout \\
Delete row & --- & Layout \\
Insert column & --- & Layout \\
Delete column & --- & Layout \\
Bold & \key{Ctrl+B} & Layout \\
Column width & --- & Layout \\
Freeze & --- & Layout \\
Sort up & --- & Data \\
Sort down & --- & Data \\
Undo sort & --- & Data \\
Recalc now & --- & Data \\
Auto recalc & --- & Data \\
Bars & --- & Chart \\
Line & --- & Chart \\
Pie & --- & Chart \\
Next & --- & Sheet \\
Previous & --- & Sheet \\
Add\ccdots & --- & Sheet \\
Rename\ccdots & --- & Sheet \\
Delete & --- & Sheet \\
\end{cctablethree}

\section{Alphabetical reference}

\cccmd[Sheet menu. No keyboard shortcut.]{Add\ccdots}
\index{Add command}

Adds a new, empty worksheet to the end of the workbook.

\ccrunin{Procedure}
\begin{ccprocedure}[]
\item Choose \menupath{Sheet > Add\ccdots}
\item Type a name of up to thirty-one characters.
\item Press \key{Return}.
\end{ccprocedure}

\ccrunin{Notes}
\begin{ccsteps}
\item The new worksheet is \emph{not} switched to. You stay where you were.
\item A tab shows only the first nine characters of a name.
\item An empty name cancels the command.
\item The seventeenth worksheet fails with \msg{Resource limit reached}.
\end{ccsteps}

\cccmd[Data menu. No keyboard shortcut.]{Auto recalc}
\index{Auto recalc command}

Turns automatic recalculation off, or back on again. It is one command and a
switch, not two commands.

\ccrunin{Notes}
\begin{ccsteps}
\item \ccprod{} starts with automatic recalculation on.
\item With it off, the right-hand end of the status line reads \msg{CALC} once
something has changed and the answers on screen have gone stale. Nothing is
shown while they are still current.
\item Turning it back on recalculates at once if anything was skipped. Turning
it off does nothing immediately.
\item See Chapter~\ref{ch:recalc}.
\end{ccsteps}

\cccmd[Chart menu. No keyboard shortcut.]{Bars}
\index{Bars command}\index{chart!bar}

Draws the numbers in the cursor's column, from the cursor's row down, as a bar
chart filling the screen. \key{S} saves it as \fname{CHART.BMP}; any other key
returns to the worksheet.

\ccrunin{Notes}
\begin{ccsteps}
\item At most sixteen values. Text, booleans and empty cells are skipped.
\item Labels come from the column immediately to the left.
\item The scale always includes zero; negative values hang below the zero line.
\item \msg{No numbers} if there is nothing to draw.
\item See Chapter~\ref{ch:charting}.
\end{ccsteps}

\cccmd[Edit menu. \key{Backspace} or \key{Delete}.]{Clear}
\index{Clear command}

Empties the current cell --- its value, its formula and its format.

\ccrunin{Notes}
\begin{ccsteps}
\item Acts on the marked block if there is one, and on the current cell if not.
\item A single cleared cell can be undone with \key{F5}, once. A cleared block
cannot.
\item A cleared cell is empty, not zero. \fn{SUM} and \fn{COUNT} skip it; a
formula reading it directly gets 0.
\end{ccsteps}

\cccmd[\key{Ctrl+B}, or the Layout menu.]{Bold}
\index{Bold command}\index{bold}

Draws the selected cells in a heavier face. One command for both directions:
if the first cell of the block is already bold the whole block comes off,
otherwise it goes on --- so a second use undoes the first.

With no block selected it acts on the cell the cursor is in.

Bold is part of the cell's format, not its contents. It is saved in an
\fname{.X16S}, written into an \fname{.xlsx} that Excel opens as bold, and
read back from one.

\begin{ccnote}
The bold face is built when \ccprod{} starts, by thickening the machine's own
character set into the upper half of the character codes. Those codes hold
accented letters, which \ccprod{} does not use --- so bold costs a second
face rather than any memory, and there is only room for the one.
\end{ccnote}

\begin{cccaution}
\key{Ctrl+B} and \key{Page Down} send the machine the same code, \lit{\$02}.
\ccprod{} tells them apart by asking the KERNAL which modifier keys are held
at that moment, so hold \key{Ctrl} until \key{B} is down. Releasing it in the
same instant pages down instead.
\end{cccaution}

\cccmd[Layout menu. No keyboard shortcut.]{Column width}
\index{Column width command}\index{columns!width}

Sets the width of the cursor's column, in characters.

\ccrunin{Procedure}
\begin{ccprocedure}[]
\item Put the cursor anywhere in the column.
\item Choose \menupath{Layout > Column width}.
\item Type a number and press \key{Return}. \key{Esc} cancels.
\end{ccprocedure}

\ccrunin{Notes}
\begin{ccsteps}
\item From 3 to 40 characters. The default is 10.
\item A worksheet holds at most \textbf{nine} widths that differ from the
default. Setting a tenth does nothing.
\item Widths are saved with the workbook, and are written to and read from
\fname{.XLSX}.
\end{ccsteps}

\cccmd[Edit menu. \keyc{Ctrl}{C}.]{Copy}
\index{Copy command}

Copies the marked block to the clipboard, or the current cell if no block is
marked. \ccprod{} says \msg{Copied}.

\ccrunin{Notes}
\begin{ccsteps}
\item What is copied is each cell's \emph{contents} as the formula bar shows
them: the formula, or the unformatted number, or the text.
\item The number format is not copied.
\item The clipboard survives switching worksheets, so a block can be copied
from one sheet and pasted into another. It does not survive quitting.
\end{ccsteps}

\cccmd[Sheet menu. No keyboard shortcut.]{Delete}
\index{Delete command (sheet)}

Deletes the current worksheet and every cell on it.

\ccrunin{Procedure}
\begin{ccprocedure}[]
\item Go to the worksheet you want to delete.
\item Choose \menupath{Sheet > Delete}.
\item Press \key{Y} at \msg{delete it and its cells?}
\end{ccprocedure}

\ccrunin{Notes}
\begin{ccsteps}
\item \textbf{Cannot be undone.}
\item The last remaining worksheet cannot be deleted.
\item Formulas on other worksheets that referred to this one keep their last
calculated values.
\end{ccsteps}

\cccmd[Layout menu. No keyboard shortcut.]{Delete column}
\index{Delete column command}\index{columns!deleting}

Removes the cursor's column, pulling everything to its right one column left.

\ccrunin{Notes}
\begin{ccsteps}
\item Asks \msg{Delete column? (y/n)} first. \key{Y} goes ahead; anything else
stops.
\item \textbf{There is no undo.}
\item Formula references are rewritten so that they go on meaning the same
cells.
\item Acts on the current worksheet only.
\end{ccsteps}

\cccmd[Layout menu. No keyboard shortcut.]{Delete row}
\index{Delete row command}

Removes the cursor's row from the current worksheet and moves everything below
it up one.

\ccrunin{Procedure}
\begin{ccprocedure}[]
\item Put the cursor anywhere on the row.
\item Choose \menupath{Layout > Delete row}.
\item Press \key{Y} at \msg{Delete row? (y/n)} on the bottom line.
\end{ccprocedure}

\ccrunin{Notes}
\begin{ccsteps}
\item \textbf{Cannot be undone.}
\item \textbf{Formula references are adjusted} --- except those naming another
worksheet, those pointing straight at the deleted row, and formulas over 128
characters long. See Chapter~\ref{ch:rows}.
\item Does nothing if the worksheet is empty or the cursor is below the last
used row.
\end{ccsteps}

\cccmd[Edit menu. \key{F2}.]{Edit cell}
\index{Edit cell command}

Puts the current cell's contents on the formula bar for editing, rather than
replacing them.

\ccrunin{Notes}
\begin{ccsteps}
\item The insertion point starts at the end of the line.
\item What appears is the contents, not the display: a formula shows as a
formula, and a number shows unformatted.
\item \key{Esc} abandons the edit and leaves the cell untouched.
\end{ccsteps}

\cccmd[File menu. \key{F8}.]{Export}
\index{Export command}

Writes the workbook out in a foreign format.

\ccrunin{Procedure}
\begin{ccprocedure}[]
\item Choose \menupath{File > Export} or press \key{F8}.
\item Type a file name.
\item Press \key{Return}.
\end{ccprocedure}

\ccrunin{Notes}
\begin{ccsteps}
\item A name ending \fname{.XLSX}, in any case, writes an Excel workbook of
every worksheet. Anything else writes a comma-separated file of the
\emph{current} worksheet only.
\item No extension is added for you.
\item An exported \fname{.XLSX} is uncompressed and about five times the size
Excel would write.
\item A CSV export writes displayed values, not formulas.
\item Nothing is shown on success.
\end{ccsteps}

\cccmd[Edit menu. \keyc{Ctrl}{F}.]{Find}
\index{Find command}

Moves the cursor to the next cell on the current worksheet whose displayed text
contains what you type.

\ccrunin{Notes}
\begin{ccsteps}
\item A substring match, ignoring case, against the cell's \emph{displayed}
text --- its result, not its formula --- and only the first 31 characters of it.
\item Searching starts after the cursor and wraps round, so repeating the
command walks through the matches. The search text is remembered between
searches, which makes \keyc{Ctrl}{F} \key{Return} a find-next.
\item At most fifteen characters, typed forwards, deleted with
\key{Backspace}. No insertion point.
\item Current worksheet only.
\item \textbf{A failed search says nothing} and leaves the cursor alone.
\end{ccsteps}

\cccmd[Layout menu. No keyboard shortcut.]{Freeze}
\index{Freeze command}

Pins the rows above the cursor and the columns to its left so that they stay on
screen while the rest scrolls. Choosing it again unfreezes.

\ccrunin{Procedure}
\begin{ccprocedure}[]
\item Put the cursor on the first cell of the data --- \cellref{B2} if the
headings are in row~1 and column~\cellref{A}.
\item Choose \menupath{Layout > Freeze}.
\end{ccprocedure}

\ccrunin{Notes}
\begin{ccsteps}
\item At \cellref{A1}, freezes the first row and the first column.
\item Rows are frozen only if the cursor is in the top half of the visible
rows; columns only if another column of default width would still fit across
the screen. Otherwise that half of the request is refused, silently.
\item The freeze is not saved with the workbook.
\end{ccsteps}

\cccmd[File menu. \key{F7}.]{Import}
\index{Import command}

Reads a foreign file into the workbook.

\ccrunin{Procedure}
\begin{ccprocedure}[]
\item For a CSV, first put the cursor where you want the top-left corner of the
data to land.
\item Choose \menupath{File > Import} or press \key{F7}.
\item Choose a file from the list and press \key{Return}.
\end{ccprocedure}

\ccrunin{Notes}
\begin{ccsteps}
\item \textbf{A CSV is read in at the cursor. An XLSX replaces the whole
workbook.} Save first.
\item The list shows only \fname{.XLSX} and \fname{.CSV} files, and only the
first sixty-four entries.
\item The dialogue is titled \msg{Open}.
\item Formulas an XLSX contains that \ccprod{} cannot compile are kept as their
last calculated values and will not update. See Chapter~\ref{ch:importexport}.
\end{ccsteps}

\cccmd[Layout menu. No keyboard shortcut.]{Insert column}
\index{Insert column command}\index{columns!inserting}

Opens a blank column at the cursor, pushing everything from there rightwards
one column along.

\ccrunin{Notes}
\begin{ccsteps}
\item Formula references are rewritten so that they go on meaning the same
cells.
\item Does nothing if the last column in use is \cellref{IU} or later --- there
would be nowhere for the last column to go.
\item Acts on the current worksheet only. Cannot be undone; delete the column
again to reverse it.
\end{ccsteps}

\cccmd[Layout menu. No keyboard shortcut.]{Insert row}
\index{Insert row command}

Opens a blank row above the cursor's row.

\ccrunin{Notes}
\begin{ccsteps}
\item No confirmation is asked.
\item \textbf{Formula references are adjusted}, with the three exceptions in
Chapter~\ref{ch:rows}.
\item Does nothing if the worksheet is empty or the cursor is below the last
used row; fails if the last used row is 65,534.
\end{ccsteps}

\cccmd[Chart menu. No keyboard shortcut.]{Line}
\index{Line command}\index{chart!line}

The same values as \menupath{Bars}, against the same scale, drawn as a red
trace with a thicker mark at each value instead of as blue blocks. \key{S}
saves it as \fname{CHART.BMP}.

\cccmd[File menu. \key{F1}.]{New}
\index{New command}

Discards the workbook in memory and starts an empty one.

\ccrunin{Notes}
\begin{ccsteps}
\item Asks \msg{discard unsaved changes?} if the workbook has been modified.
\item Gives one worksheet, the cursor on \cellref{A1}, and no file name.
\item A failure here is reported as \msg{Save failed}, which is a slip in the
program rather than anything to do with saving.
\end{ccsteps}

\cccmd[Sheet menu. No keyboard shortcut.]{Next}
\index{Next command}

Goes to the next worksheet, wrapping round from the last to the first.

\ccrunin{Notes}
\begin{ccsteps}
\item The cursor is put on \cellref{A1}. Your position on the old worksheet is
not remembered.
\end{ccsteps}

\cccmd[File menu. \key{F3}.]{Open}
\index{Open command}

Loads a \fname{.X16S} workbook.

\ccrunin{Procedure}
\begin{ccprocedure}[]
\item Choose \menupath{File > Open} or press \key{F3}.
\item Answer \msg{discard unsaved changes?} if it is asked.
\item Move the highlight with \key{$\uparrow$} and \key{$\downarrow$}.
\item Press \key{Return}.
\end{ccprocedure}

\ccrunin{Notes}
\begin{ccsteps}
\item The list holds \textbf{at most sixty-four} entries, sorted, sixteen shown
at a time, with \lit{\^{}} and \lit{v} marking more above and below.
\item Subdirectories are listed first, as \lit{/NAME}; \lit{/..} goes back up.
Selecting one and pressing \key{Return} lists it. The path is shown on the
title row.
\item Browsing is not remembered --- the next \menupath{Open} starts at the
directory \ccprod{} was started in.
\item Only \fname{.X16S} files are listed.
\item An invalid file is rejected before the workbook in memory is touched. A
file that fails part way through loading leaves an empty workbook.
\end{ccsteps}

\cccmd[Edit menu. \keyc{Ctrl}{V}.]{Paste}
\index{Paste command}

Enters the clipboard's contents into the current cell.

\ccrunin{Notes}
\begin{ccsteps}
\item The cursor is the \textbf{top-left corner} of what arrives.
\item Exactly as though you had typed them --- so a formula is compiled again,
and a number arrives with no format.
\item \textbf{A pasted formula is rewritten for its new position}, by the
distance the block moved, in rows and columns at once. \fml{=A1+B1} copied five
rows down and two across becomes \fml{=C6+D6}. A \lit{\$} pins what should not
move.
\item Not undoable. \key{F5} restores one cell, and a block is not one cell.
\end{ccsteps}

\cccmd[Chart menu. No keyboard shortcut.]{Pie}
\index{Pie command}\index{chart!pie}

Draws the same values as a circle divided into slices, with a key beside it
giving each slice's label and its share as a whole-number percentage. \key{S}
saves it as \fname{CHART.BMP}.

\ccrunin{Notes}
\begin{ccsteps}
\item \textbf{Negative values are dropped}, and the percentages are shares of
the total of what is left.
\item If nothing is left, \msg{No numbers}.
\item Percentages are whole numbers rounded down, so they need not total 100.
\item Slice colours cycle through six; with more than six values, two slices
share a colour and the key is what tells them apart.
\end{ccsteps}

\cccmd[Sheet menu. No keyboard shortcut.]{Previous}
\index{Previous command}

Goes to the previous worksheet, wrapping round from the first to the last. The
cursor is put on \cellref{A1}.

\cccmd[File menu. \key{F9}.]{Quit}
\index{Quit command}

Leaves \ccprod{} and returns to BASIC.

\ccrunin{Notes}
\begin{ccsteps}
\item Asks \msg{discard unsaved changes?} if the workbook has been modified.
\item \key{Y} quits. Any other key cancels.
\end{ccsteps}

\cccmd[Data menu. No keyboard shortcut.]{Recalc now}
\index{Recalc now command}

Recalculates every formula in the workbook immediately, whether automatic
recalculation is on or off.

\ccrunin{Notes}
\begin{ccsteps}
\item This is what clears \msg{CALC} from the status line.
\item With automatic recalculation on it has nothing to do, but it is harmless.
\end{ccsteps}

\cccmd[Sheet menu. No keyboard shortcut.]{Rename\ccdots}
\index{Rename command}

Changes the name of the current worksheet. The existing name appears in the
field ready to be edited.

\ccrunin{Notes}
\begin{ccsteps}
\item Formulas on other worksheets that name this one by its old name are not
rewritten and are not re-checked. They keep their last values.
\end{ccsteps}

\cccmd[Edit menu. No keyboard shortcut.]{Replace\ccdots}
\index{Replace command}

Finds text throughout the worksheet and replaces it, asking at each match.

\ccrunin{Procedure}
\begin{ccprocedure}[]
\item Choose \menupath{Edit > Replace\ccdots}.
\item Type what to find, and press \key{Return}.
\item Type what to put in its place, and press \key{Return}. An empty
replacement deletes the matched text.
\item At each match the cursor moves there and \ccprod{} asks
\msg{Replace? Y/N/A, Esc to stop}.
\end{ccprocedure}

\ccrunin{Notes}
\begin{ccsteps}
\item \key{Y} replaces this one, \key{N} skips it, \key{A} replaces every
remaining match without asking, \key{Esc} stops.
\item \textbf{The whole used range, from \cellref{A1}} --- not from the cursor,
unlike \menupath{Find}.
\item \textbf{Formula cells are never touched.}
\item Every occurrence within a cell is replaced, and matching ignores case.
\item Numbers are matched as well as text, and a replacement can change a cell
from one to the other.
\item \textbf{Cannot be undone.} Save first.
\item See Chapter~\ref{ch:finding}.
\end{ccsteps}

\cccmd[File menu. \key{F4}.]{Save}
\index{Save command}

Writes the workbook in \ccprod's own \fname{.X16S} format.

\ccrunin{Notes}
\begin{ccsteps}
\item If the workbook already has a name it is saved under that name with no
dialogue. Otherwise you are asked for one.
\item \textbf{Always writes \fname{.X16S}}, whatever extension the name has.
\item The save is atomic: written to \fname{CMDRCALC.TMP}, verified, and only
then renamed over the old file.
\item Needs room on the card for two copies of the workbook.
\item Nothing is shown on success; the status line changes from \msg{MODIF}
back to \msg{READY}.
\end{ccsteps}

\cccmd[File menu. \key{F6}.]{Save as}
\index{Save as command}

The same as \menupath{Save}, except that it always asks for a name.

\cccmd[Edit menu. \keyc{Ctrl}{A}.]{Select}
\index{Select command}\index{block!selecting}

Anchors one corner of a block at the cursor, or drops the block if one is
already marked.

\ccrunin{Notes}
\begin{ccsteps}
\item Once anchored, \textbf{every movement key extends the block} --- arrows,
\key{PgUp}, \key{PgDn}, \key{Home}, \key{End}, and a mouse click.
\item \keyc{Ctrl}{A} again drops it. \key{Esc} does not --- it belongs to
editing a cell.
\item \menupath{Copy} and \menupath{Clear} drop the block once they are done
with it. Pasting and sorting leave it marked.
\item The anchor may end up below and right of the cursor; the block is the
rectangle between them either way.
\item \menupath{Copy}, \menupath{Clear} and the two \menupath{Sort} commands
act on the block when there is one and on the cursor cell or column when there
is not.
\end{ccsteps}

\cccmd[Data menu. No keyboard shortcut.]{Sort down}
\index{Sort down command}

Sorts rows into descending order on the cursor's column, from the cursor's row
down. See \menupath{Sort up}; everything there applies, and empty rows still
sort to the bottom.

\cccmd[Data menu. No keyboard shortcut.]{Sort up}
\index{Sort up command}

Sorts into ascending order. With a block selected, only that block moves,
ordered on its own first column. With nothing selected, whole rows are sorted
on the column the cursor is in, from the cursor's row down.

\ccrunin{Procedure}
\begin{ccprocedure}[]
\item Either select the table --- \keyc{Ctrl}{A} and extend --- or put the
cursor in the column to sort by, on the first row of data, not on the heading.
\item Choose \menupath{Data > Sort up}.
\item If the range contains formulas, press \key{Y} at
\msg{ranges over sorted rows may be wrong}.
\item If the block is narrower than the data on its rows, press \key{Y} at
\msg{data beside the block will not move}.
\end{ccprocedure}

\ccrunin{Notes}
\begin{ccsteps}
\item Order: numbers, then text (case-insensitive, first thirty-two
characters), then errors, then blanks. Blanks stay last in both directions.
\item \textbf{Relative references move with their row}, so a formula reading
its own row goes on reading it. A \emph{range} spanning sorted rows cannot be
fixed that way, which is what the warning is about.
\item At most 4,096 rows; beyond that, \msg{too many rows to sort}.
\item \textbf{A selection sorts on its own first column}, not on the sheet's
and not on the cursor's: a block from \cellref{C1} to \cellref{E9} orders by
column~\cellref{C}.
\item The cursor is put on \cellref{A1} afterwards, and a selection is left
marked.
\end{ccsteps}

\cccmd[Edit menu. \key{F5}.]{Undo}
\index{Undo command}

Restores the last single cell that was changed.

\ccrunin{Notes}
\begin{ccsteps}
\item One cell, one level. \textbf{Undoing is not itself undoable.}
\item Does not reach a pasted or cleared block, a sort, an inserted or deleted
row or column, a deleted worksheet, an import, or an opened file.
\end{ccsteps}

\cccmd[Data menu. No keyboard shortcut.]{Undo sort}
\index{Undo sort command}

Puts the rows back the way they were before the last sort.

\ccrunin{Notes}
\begin{ccsteps}
\item Works once, then not again.
\item Refuses, silently, if you are on a different worksheet or the used range
has changed size.
\item Forgotten by any file command.
\item Unrelated to \key{F5} \menupath{Undo}.
\end{ccsteps}
